El presente informe trata el análisis experimental de algoritmos aplicados a dos problemas fundamentales: el ordenamiento de arreglos unidimensionales de números enteros y la multiplicación de matrices cuadradas. Si bien sus complejidades teóricas son conocidas, el objetivo principal de este trabajo es evaluar su comportamiento empírico en la práctica. Mediante la medición de tiempos de ejecución y consumo de memoria en distintos escenarios, se busca relacionar directamente el rendimiento práctico de las implementaciones con su complejidad teórica asintótica esperada.


